\documentclass[10pt,
  aspectratio=169,
  serif,
  mathserif,
  professionalfont,
  compress,
  handout,
  % table,
  % svgnames
  ]{beamer}\usepackage[]{graphicx}\usepackage[]{color}
% maxwidth is the original width if it is less than linewidth
% otherwise use linewidth (to make sure the graphics do not exceed the margin)
\makeatletter
\def\maxwidth{ %
  \ifdim\Gin@nat@width>\linewidth
    \linewidth
  \else
    \Gin@nat@width
  \fi
}
\makeatother

\definecolor{fgcolor}{rgb}{1, 1, 0.941}
\newcommand{\hlnum}[1]{\textcolor[rgb]{0.804,0.718,0.71}{#1}}%
\newcommand{\hlstr}[1]{\textcolor[rgb]{0.604,0.753,0.804}{#1}}%
\newcommand{\hlcom}[1]{\textcolor[rgb]{0.439,0.502,0.565}{#1}}%
\newcommand{\hlopt}[1]{\textcolor[rgb]{1,1,0.941}{#1}}%
\newcommand{\hlstd}[1]{\textcolor[rgb]{1,1,0.941}{#1}}%
\newcommand{\hlkwa}[1]{\textcolor[rgb]{0.941,0.902,0.549}{#1}}%
\newcommand{\hlkwb}[1]{\textcolor[rgb]{1,0.871,0.678}{#1}}%
\newcommand{\hlkwc}[1]{\textcolor[rgb]{0.545,0.941,0.702}{#1}}%
\newcommand{\hlkwd}[1]{\textcolor[rgb]{0.545,0.941,0.902}{#1}}%
\let\hlipl\hlkwb

\usepackage{framed}
\makeatletter
\newenvironment{kframe}{%
 \def\at@end@of@kframe{}%
 \ifinner\ifhmode%
  \def\at@end@of@kframe{\end{minipage}}%
  \begin{minipage}{\columnwidth}%
 \fi\fi%
 \def\FrameCommand##1{\hskip\@totalleftmargin \hskip-\fboxsep
 \colorbox{shadecolor}{##1}\hskip-\fboxsep
     % There is no \\@totalrightmargin, so:
     \hskip-\linewidth \hskip-\@totalleftmargin \hskip\columnwidth}%
 \MakeFramed {\advance\hsize-\width
   \@totalleftmargin\z@ \linewidth\hsize
   \@setminipage}}%
 {\par\unskip\endMakeFramed%
 \at@end@of@kframe}
\makeatother

\definecolor{shadecolor}{rgb}{.97, .97, .97}
\definecolor{messagecolor}{rgb}{0, 0, 0}
\definecolor{warningcolor}{rgb}{1, 0, 1}
\definecolor{errorcolor}{rgb}{1, 0, 0}
\newenvironment{knitrout}{}{} % an empty environment to be redefined in TeX

\usepackage{alltt}

% Tamanho de fonte e distância entre linhas.
\renewenvironment{knitrout}{
  \renewcommand{\baselinestretch}{0.75}%\tiny
}{}

%-----------------------------------------------------------------------
% Pacotes padrões.

% Fontes.
\usepackage{palatino}
\usepackage{eulervm}
\usepackage{inconsolata}

% Esses pacotes dão clash.
% http://tex.stackexchange.com/questions/51488/option-clash-with-xcolor-and-tikz
% \usepackage{xcolor} %% opções no \documentclass{} para evitar clash.
% \definecolor{mycolor}{rgb}{0.13,0.53,0.53}
% \definecolor{mycolor2}{rgb}{0.725,0,0.18}

\usepackage{hyperref}
 \hypersetup{colorlinks, allcolors=., urlcolor=structure}
% \hypersetup{colorlinks}

\usepackage[brazil]{babel}
\usepackage[utf8]{inputenc}
\usepackage{graphicx}
\usepackage{amsmath, amsfonts, amssymb, amsxtra, amsthm, icomma}
\usepackage{geometry, calc, setspace, indentfirst}
% \usepackage{colortbl}
\usepackage{enumerate}
\usepackage{float}

\usepackage{subfigure}

\usepackage[hang]{caption}
\captionsetup{font=footnotesize,
  labelfont=footnotesize,
  labelsep=period}

% Listas em duas colulas.
\usepackage{multicol}
\newenvironment{itemize2}{%
  \vspace*{-1em}
  \begin{itemize}
    \begin{multicols}{2}
    }{%
    \end{multicols}
  \end{itemize}
}

% Texto no corpo do beamer justificado.
\usepackage{ragged2e}
\justifying

%-----------------------------------------------------------------------

% A lot of options:
% http://latex-community.org/forum/viewtopic.php?f=55&t=17646
\useoutertheme[
  width=60pt,
  height=30pt,
  right,
  hideothersubsections
  ]{sidebar}

\makeatletter
\setbeamertemplate{caption}[numbered]
\setbeamertemplate{section in toc}[sections numbered]
\setbeamertemplate{subsection in toc}[subsections numbered]
\setbeamertemplate{sections/subsections in toc}[ball]{}
\setbeamertemplate{section in sidebar right}[sections numbered]
%\setbeamertemplate{frametitle continuation}{\gdef\beamer@frametitle{}}
\setbeamertemplate{navigation symbols}{} %% Retira a barra de navegação.
\setbeamertemplate{footline}[frame number]
% \setbeamertemplate{blocks}[rounded][shadow=FALSE]
% \setbeamercolor{block title}{fg=structure, bg=mycolor!20!white}
\makeatother

% Frames com sessão e/ou subsessão.
% \addtobeamertemplate{frametitle}{
%   \let\insertframetitle\insertsubsectionhead}{}
% \makeatletter
% \CheckCommand*\beamer@checkframetitle{
%   \@ifnextchar\bgroup\beamer@inlineframetitle{}}
% \renewcommand*\beamer@checkframetitle{
%   \global\let\beamer@frametitle\relax\@ifnextchar\bgroup
%   \beamer@inlineframetitle{}}
% \makeatother

\setbeamertemplate{frametitle}
{
    \nointerlineskip
    \begin{beamercolorbox}[sep=0.3cm,ht=1.8em,wd=\paperwidth]{frametitle}
        \vbox{}\vskip-2ex%
        \strut\insertframetitle\strut
        \vskip-0.8ex%
    \end{beamercolorbox}
}


%-----------------------------------------------------------------------
% Comandos.

\newcommand{\n}[1]{\textbf{#1}}

%-----------------------------------------------------------------------

\AtBeginSection[]{
  \begin{frame}[c,allowframebreaks]
    \begin{center}
      {\thesection} \\ \vspace{0.3cm}
      \parbox{0.6\textwidth}{
        \centering {\Large \textcolor{structure}{\insertsection}}}\\
    \end{center}
  \end{frame}
}

%-----------------------------------------------------------------------
% Definições dos proprietários.

\title[TH MCGLM]{
  \Large Testes de hipóteses em Modelos Multivariados de \\ Covariância Linear Generalizada\\ \vspace{0.3cm}}

\subtitle{Semana acadêmica do curso de Estatística 2021}

\author[Lineu Alberto]{\small
  Lineu Alberto Cavazani de Freitas \\
  Orientador: Prof. Dr. Wagner Hugo Bonat\\
  Co-orientador: Prof. Dr. Marco Antonio Zanata Alves
}

 \institute[UFPR]{
    PPG Informática %\\
%   Data Science \& Big Data \\
%   Universidade Federal do Paraná\\
% 
%   \vspace{1em}
%   \href{}{https://lineu96.github.io/st/}\\
%   \texttt{lineuacf@gmail.com}
 }
 \date{}

\logo{\includegraphics[width=1.5cm]{img/dsbd1x4-rect.png}} 

\usebackgroundtemplate{
  \includegraphics[width=\paperwidth]{img/ufpr-fundo.jpg}
}

\titlegraphic{
  \vspace{-1em}
  \includegraphics[height=1.8cm]{img/capes_tp2.png}\hspace{2em}
  \includegraphics[height=1.8cm]{img/ufpr-transparent.png}\hspace{2em}
  \includegraphics[height=1.8cm]{img/dsbd-2x2-trans.png}
}

%---- preamble-refs.tex ------------------------------------------------

% Bibliography.

% %\usepackage[style=authoryear]{biblatex}
% \usepackage[authordate, bibencoding=auto, strict, backend=biber, natbib]{biblatex-chicago}
%
% % Use:
% %   \cite{<ref>}
% %   \parencite{<ref>}
% %   \fullcite{<ref>}
% %   \footfullcite{<ref>}
%
% % ATTENTION
% % Compilation: pdflatex > biber > pdflatex > pdflatex.
%
% % Calls refs.bib at preamble with:
% \addbibresource{config/refs.bib}

% abntex2cite -------------------------------

\usepackage[
  alf,
  abnt-emphasize=bf,
  abnt-etal-list=2,
  abnt-and-type=&]{abntex2cite}
% Use:
%   \cite{<ref>}
%   \citeonline{<ref>}

% ATTENTION
% Compilation: pdflatex > bibtex > pdflatex > pdflatex.

% Calls refs.bib at last frame with:
% \bibliography{config.refs}

\let\oldbibliography\thebibliography
\renewcommand{\thebibliography}[1]{%
  \oldbibliography{#1}%
  \setlength{\itemsep}{1em}%
}

%--------------------------------------------

\usepackage{pgfgantt}

%=======================================================================
%=======================================================================
\IfFileExists{upquote.sty}{\usepackage{upquote}}{}
\begin{document}

\frame{\titlepage}

% Tabela de conteúdo no início dos slides.
 \begin{frame}{Conteúdo}
   \small{\tableofcontents}
 \end{frame}

%-----------------------------------------------------------------------



%-----------------------------------------------------------------------

\section{Ciência de dados}

%-----------------------------------------------------------------------

\begin{frame}
  \frametitle{Ciência de dados}
  \begin{itemize}
    \itemsep 2ex
  
  \item \textbf{Ciência de dados} é campo de estudo interdisciplinar que incorpora conhecimento de áreas como:
  
  \begin{enumerate}
    \item \textbf{Estatística}.
    \item \textbf{Ciência da computação}.
    \item \textbf{Matemática}.
  \end{enumerate}

  \item Os \textbf{métodos estatísticos} são de fundamental importância em grande parte das etapas da ciência de dados \cite{weihs2018data}.
  
  \item Neste sentido, os \textbf{modelos de regressão} tem papel importante.
  
  \end{itemize}
\end{frame}

%-----------------------------------------------------------------------

\section{Modelos de regressão}

%-----------------------------------------------------------------------

\begin{frame}
  \frametitle{Modelos de regressão}
  
  \begin{itemize}
  
  \itemsep 2ex
  
  \item Técnica utilizada quando há o interesse em verificar a influência das variáveis explicativas sobre variáveis respostas.
  
  \item Assume-se que a variável resposta condicional às explicativas segue uma distribuição de probabilidades. 
  
  \item Um modelo de regressão é uma expressão matemática que relaciona a média da variável resposta às variáveis explicativas.

  \item Com modelos de regressão é possível:
  
  \begin{itemize}
    \item \textbf{Explicar} a associação entre as variáveis.
    \item \textbf{Predizer} valores da resposta com base nos valores observados das explicativas. 
  \end{itemize}

\end{itemize}

\end{frame}

%-----------------------------------------------------------------------


\begin{frame}
  \frametitle{Modelos de regressão}

  \begin{enumerate}
    \itemsep 2ex
    
      \item \textbf{Definição do problema.}
      
      \begin{itemize}
        \item Qual o fenômeno aleatório de interesse?
        \item Que fatores externos podem afetar este fenômeno?
      \end{itemize}

      
      \item \textbf{Planejamento do estudo e coleta de dados.}
      
      \item \textbf{Análise dos dados via regressão.}
      
        \begin{itemize}
          \item Escolha da distribuição de probabilidade.
          \item Especificação do modelo.
          \item Obtenção das estimativas dos parâmetros (ajuste).
          \item Diagnóstico.
        \end{itemize}
      
      \item \textbf{Interpretação dos resultados.}
      
        \begin{itemize}
          \item Quais os fatores externos apresentam ou não impacto sobre o fenômeno.
          \item Qual a dimensão desse impacto.
        \end{itemize}      
  
  \end{enumerate}
\end{frame}

%-----------------------------------------------------------------------
\section{Modelos multivariados de covariância linear generalizada}
%-----------------------------------------------------------------------

\begin{frame}

\frametitle{Modelos multivariados de covariância linear generalizada}

\begin{itemize}
   
   \itemsep 2ex
   
  \item Os modelos multivariados de covariância linear generalizada (McGLMs) \cite{Bonat16} configuram uma estrutura geral para análise via modelos de regressão.
  
  \item Comporta múltiplas respostas de diferentes naturezas.
  
    \item Pode-se ajustar modelos com diferentes preditores e distribuições para cada resposta.
  
  \item Pode especificar modelos que levam em conta a correlação entre indivíduos do conjunto de dados. 

\end{itemize}

\end{frame}

%-----------------------------------------------------------------------

\begin{frame}

\frametitle{Modelos multivariados de covariância linear generalizada}

\begin{itemize}
   
   \itemsep 2ex
  
  \item Os parâmetros são interpretáveis:
  
    \begin{itemize}
    
    \itemsep 2ex
      
      \item \textbf{Parâmetros de regressão}: efeito das variáveis explicativas sobre as respostas.
      
      \item \textbf{Parâmetros de dispersão}: impacto da correlação entre unidades.
      
      \item \textbf{Parâmetros de potência}: indicativo de qual distribuição se adequa ao problema.
      
    \end{itemize}
    
    \item A metodologia do McGLM está implementada no pacote \emph{mcglm} \cite{mcglm} do software R.

\end{itemize}

\end{frame}

%-----------------------------------------------------------------------

\begin{frame}
\frametitle{Modelos multivariados de covariância linear generalizada}

Para definição de um McGLM considere:

\begin{itemize}
  
  \itemsep 2ex
  
  \item $\boldsymbol{Y}_{N \times R} = \left \{ \boldsymbol{Y}_1, \dots, \boldsymbol{Y}_R \right \}$ uma  matriz de variáveis resposta.
  
  \item $\boldsymbol{M}_{N \times R} = \left \{ \boldsymbol{\mu}_1, \dots, \boldsymbol{\mu}_R \right \}$ uma matriz de valores esperados.
  
    \item $\boldsymbol{X}_r$ denota uma matriz de delineamento $N \times k_r$.
  
  \item $\boldsymbol{\beta}_r$ denota um vetor $k_r \times 1$ de parâmetros de regressão.
  
\end{itemize}

\end{frame}

%-----------------------------------------------------------------------

\begin{frame}
\frametitle{Modelos multivariados de covariância linear generalizada}

Considere ainda:

\begin{itemize}
  
  \itemsep 2ex
  
    \item $\Sigma_b$ uma matriz de correlação entre variáveis resposta, de ordem $R \times R$.
    
  \item $\Sigma_r$, $r = 1,..., R$, a matriz de variância e covariância para cada resposta $r$, de dimensão $NxN$:
  
$$
\Sigma_r = \mathrm{V}_r\left(\boldsymbol{\mu}_r; p_r\right)^{1/2}(\boldsymbol{\Omega}\left(\boldsymbol{\tau}_r\right))\mathrm{V}_r\left(\boldsymbol{\mu}_r; p_r\right)^{1/2}.
$$

Em que:

  \begin{itemize}
    \item $\mathrm{V}_r\left(\boldsymbol{\mu}; p\right)$ é uma matriz diagonal em que as entradas principais são dadas pela função de variância aplicada ao vetor $\boldsymbol{\mu}$. 
  
  \item $p_r$ é o parâmetro de potência. 
  
  \item $\boldsymbol{\Omega}\left(\boldsymbol{\tau}_r\right)$ a matriz de dispersão que descreve a parte da covariância dentro de cada variável resposta. 
  \end{itemize}
  

\end{itemize}

\end{frame}

%-----------------------------------------------------------------------

\begin{frame}

\frametitle{Preditor linear matricial}

\begin{itemize}
  
  \itemsep 2ex
  
  \item A matriz $\boldsymbol{\Omega}(\boldsymbol{\tau}_r)$ descreve a estrutura de correlação entre as observações da amostra.
  
  \item É modelada através de um preditor linear matricial combinado com uma função de ligação de covariância:

$$
h\left \{ \boldsymbol{\Omega}(\boldsymbol{\tau}_r) \right \} = \tau_{r0}Z_0 + \ldots + \tau_{rD}Z_D
$$
  
  \begin{itemize}
  
  \itemsep 2ex
  
  \item $h()$ é a função de ligação de covariância.
  
  \item $Z_{rd}$ com $d$ = 0,$\ldots$, D são matrizes que representam a estrutura de covariância presente em cada variável resposta $r$.
  
  \item $\boldsymbol{\tau_r}$ = $(\tau_{r0}, \ldots, \tau_{rD})$ é um vetor $(D + 1) \times 1$ de parâmetros de dispersão. 
  
\end{itemize}

     
\end{itemize}

\end{frame}

%-----------------------------------------------------------------------

\begin{frame}

  \frametitle{Funções de variância}
  
  \begin{enumerate}
  \item \textbf{Função de variância potência \cite{Jorgensen87} e \cite{Jorgensen97}. }
  
    \begin{itemize}
      \item Família Tweedie de distribuições.
      \item $\vartheta\left(\boldsymbol{\mu}; p\right) = \mu^p$.
      \item Casos particulares: normal ($p$ = 0), Poisson ($p$ = 1), gama ($p$ = 2) e normal inversa ($p$ = 3).
    \end{itemize}

  
  \item \textbf{Função de dispersão Poisson–Tweedie \cite{Jorgensen15}.}
  
    \begin{itemize}
      \item Família Poisson-Tweedie de distribuições.
      \item $\vartheta\left(\boldsymbol{\mu}; p\right) = \mu + \mu^p$.
      \item Casos particulares: Hermite ($p$ = 0), Neyman tipo A ($p$ = 1), binomial negativa ($p$ = 2) e Poisson–inversa gaussiana (p = $3$).
      
    \end{itemize}

  \item \textbf{Função de variância binomial.} 
  
    \begin{itemize}
      \item $\vartheta(\boldsymbol{\mu}) = \mu(1 - \mu)$.
      \item Acomoda respostas binárias ou restritas a um intervalo.
    \end{itemize}

\end{enumerate}

\end{frame}

%-----------------------------------------------------------------------

\begin{frame}
\frametitle{Modelos multivariados de covariância linear generalizada}

Os McGLMs são definidos por:

\begin{center}
$\mathrm{E}(\boldsymbol{Y}) = \boldsymbol{M} = \{g_1^{-1}(\boldsymbol{X}_1 \boldsymbol{\beta}_1), \ldots, g_R^{-1}(\boldsymbol{X}_R \boldsymbol{\beta}_R)\}$

$\mathrm{Var}(\boldsymbol{Y}) = \boldsymbol{C} = \boldsymbol{\Sigma}_R \overset{G} \otimes \boldsymbol{\Sigma}_b$

\end{center}
    
Em que: 

\begin{itemize}
  
  \item $\boldsymbol{\Sigma}_R \overset{G} \otimes \boldsymbol{\Sigma}_b = \mathrm{Bdiag}(\tilde{\boldsymbol{\Sigma}}_1, \ldots, \tilde{\boldsymbol{\Sigma}}_R) (\boldsymbol{\Sigma}_b \otimes \boldsymbol{I}) \mathrm{Bdiag}(\tilde{\boldsymbol{\Sigma}}_1^\top, \ldots, \tilde{\boldsymbol{\Sigma}}_R^\top)$ é o produto generalizado de Kronecker.
  
  \item $\tilde{\boldsymbol{\Sigma}}_r$ denota a matriz triangular inferior da decomposição de Cholesky da matriz ${\boldsymbol{\Sigma}}_r$.
  
  \item $\mathrm{Bdiag()}$ denota a matriz bloco-diagonal.
  
  \item $\boldsymbol{I}$ uma matriz identidade $N \times N$.
  
  \item $g_r()$ são as tradicionais funções de ligação.
  
\end{itemize}

\end{frame}

%-----------------------------------------------------------------------

\begin{frame}
\frametitle{Modelos multivariados de covariância linear generalizada}

\begin{itemize}

\itemsep 2ex
  
  \item Parâmetros estimados nos McGLMs:
    \begin{enumerate}
      \item Regressão.
      \item Dispersão.
      \item Potência.
      \item Correlação.
    \end{enumerate}
  
  \item Todas estas quantidades são interpretáveis e são estimadas com base nos dados.
  
  \item A estimação é feita por meio de \textbf{funções de estimação}.
    \begin{enumerate}
      \item \textbf{Função quasi-score} para parâmetros de regressão. 
      \item \textbf{Função de estimação de Pearson} para os demais parâmetros. 
    \end{enumerate}

  \item Mais detalhes são descritos em \cite{Bonat16}.

\end{itemize}

\end{frame}

%-----------------------------------------------------------------------

\section{Testes de hipóteses}

%-----------------------------------------------------------------------

\begin{frame}
  \frametitle{Testes de hipóteses em modelos de regressão}
  \begin{itemize}
    \itemsep 2ex
  
  \item Usados para verificar se a retirada de determinada variável explicativa do modelo geraria uma perda no ajuste.
  
  \item Dentre os testes mais comuns está o teste Wald:

    \begin{itemize}
      
      \item Verifica se existe evidência para afirmar que os parâmetros são iguais a valores postulados.
      
      \item Avalia quão longe o valor estimado está do valor postulado. 
      
        \item As hipóteses podem ser descritas como:

$$
\left\{\begin{matrix}
H_0: \boldsymbol{L}\boldsymbol{\beta} = \boldsymbol{c} \\ 
H_1: \boldsymbol{L}\boldsymbol{\beta} \neq \boldsymbol{c}
\end{matrix}\right.
$$

  \item A estatística de teste é dada por:
  $$
WT = (\boldsymbol{L\hat\beta} - \boldsymbol{c})^T \ (\boldsymbol{L \ Var^{-1}(\hat\beta) \ L^T})^{-1} \ (\boldsymbol{L\hat\beta} - \boldsymbol{c}).
$$

 \item  WT \sim \chi^2_q.
   
    \end{itemize}
  
  \end{itemize}
  
\end{frame}

%-----------------------------------------------------------------------

\begin{frame}
  \frametitle{ANOVA \& MANOVA}

  \begin{itemize}
    \itemsep 2ex

  \item Formas de \textbf{avaliar a significância} de cada uma das variáveis de uma forma procedural.  
  
  \item Consiste em efetuar testes sucessivos impondo \textbf{restrições ao modelo} original. 

  \item O objetivo é testar se a ausência de determinada variável \textbf{gera perda} ao modelo. 

  \item Os resultados são sumarizados numa tabela, o chamado \textbf{quadro de análise de variância}.
  
  \end{itemize}

\end{frame}


%-----------------------------------------------------------------------

\begin{frame}
  \frametitle{Tópicos abordados até aqui...}

\begin{enumerate}

  \itemsep 2ex
  
  \item Importância dos modelos de regressão.

  \item McGLMs e importância dos parâmetros de regressão, dispersão e potência.

  \item Testes de hipóteses em modelos de regressão, ênfase no teste Wald.
  
  \item Procedimentos baseados em testes de hipóteses: ANOVA e MANOVA.
  
\end{enumerate}

\end{frame}

%-----------------------------------------------------------------------

\section{Proposta}

%-----------------------------------------------------------------------

\begin{frame}
  \frametitle{Proposta}

\begin{itemize}

  \itemsep 2ex
  
  \item Considerando os McGLMs, não há discussão a respeito da construção de testes de hipóteses.
  
  \item Considerando o alto potencial de aplicação dos McGLMs em ciência de dados, nossos objetivos gerais são:

  \begin{itemize}
  
  \itemsep 2ex
  
  \item Desenvolvimento de testes de hipóteses para avaliação dos parâmetros de McGLMs.

  \item Adaptação do teste de Wald clássico utilizado em modelos lineares para os McGLMs. 

  \end{itemize}
  
\end{itemize}

\end{frame}

%-----------------------------------------------------------------------

\begin{frame}
  \frametitle{Etapas}

  \begin{enumerate}
    \itemsep 2ex

  \item \textbf{Adaptar} o teste Wald para realização de testes de hipóteses gerais sobre parâmetros de McGLMs. 
  
  \item \textbf{Implementar} funções para efetuar tais testes, bem como funções para efetuar ANOVAs e MANOVAs para os McGLMs. 

  \item \textbf{Avaliar} as propriedades e comportamento dos testes propostos com base em estudos de simulação.

  \item \textbf{Motivar} o potencial de aplicação das metodologias discutidas com base na aplicação a conjuntos de dados reais.

  \end{enumerate}

\end{frame}

%-----------------------------------------------------------------------

\section{Resultados preliminares}

%-----------------------------------------------------------------------

\subsection{Adaptação do teste Wald para os McGLM}

%-----------------------------------------------------------------------

\begin{frame}[c, allowframebreaks]

\begin{center}

  {\normalsize \href{https://lineu96.github.io/st/}{Adaptação do teste Wald para os McGLMs}}
  
\end{center}

\end{frame}

% -----------------------------------------------------------------

\begin{frame}
\frametitle{Hipóteses}

$$H_0: \boldsymbol{L}\boldsymbol{\theta_{\beta,\tau,p}} = \boldsymbol{c} \ vs \ H_1: \boldsymbol{L}\boldsymbol{\theta_{\beta,\tau,p}} \neq \boldsymbol{c}.$$ 

Em que: 

\begin{itemize}
  
  \item Em que $\boldsymbol{L}$ é a matriz de especificação das hipóteses a serem testadas, tem dimensão $s \times h$. 
  
  \item $\boldsymbol{\theta_{\beta,\tau,p}}$ é o vetor de dimensão $h \times 1$ de parâmetros de regressão, dispersão e potência do modelo. 
  
  \item $\boldsymbol{c}$ é um vetor de dimensão $s \times 1$ com os valores sob hipótese nula.

\end{itemize}

\end{frame}

% -----------------------------------------------------------------

\begin{frame}
\frametitle{Estatística de teste}

$$W = (\boldsymbol{L\hat\theta_{\beta,\tau,p}} - \boldsymbol{c})^T \ (\boldsymbol{L \ J_{\boldsymbol{{\beta,\tau,p}}}^{-1} \ L^T})^{-1} \ (\boldsymbol{L\hat\theta_{\beta,\tau,p}} - \boldsymbol{c}).$$

Em que: 

\begin{itemize}
  \item $\boldsymbol{L}$ é a matriz da especificação das hipóteses, tem dimensão $s \times h$. 

  \item $\boldsymbol{\hat\theta_{\beta,\tau,p}}$ é o vetor de dimensão $h \times 1$ com todas as estimativas dos parâmetros de regressão, dispersão e potência. 

  \item $\boldsymbol{c}$ é um vetor de dimensão $s \times 1$ com os valores sob hipótese nula. 

  \item E $J_{\boldsymbol{{\beta,\tau,p}}}^{-1}$ é a inversa da matriz de informação de Godambe desconsiderando os parâmetros de correlação, de dimensão $h \times h$. 

\end{itemize}

\end{frame}

% -----------------------------------------------------------------

\subsection{Exemplos de hipóteses}

% -----------------------------------------------------------------

\begin{frame}[c, allowframebreaks]

\begin{center}

  {\normalsize \href{https://lineu96.github.io/st/}{Exemplo de hipóteses nos McGLMs}}
  
\end{center}

\end{frame}

% -----------------------------------------------------------------

\begin{frame}

\frametitle{Problema}

\begin{itemize}
  
  \itemsep 2ex
  
  \item Considere o exemplo disponível na seção 4.3.6 do livro \href{https://www.ime.usp.br/~giapaula/texto_2013.pdf}{Modelos de Regressão com Apoio Computacional} \cite{paula}:
  
  \item Os dados referem-se a um estudo sobre demanda de TV's a cabo em 40 regiões dos Estados Unidos. 
  
  \item Algumas das variáveis coletadas foram:
    
    \begin{itemize}

        \item Número de assinantes de TV a cabo (em milhares).

        \item Percentual de domicílios com TV a cabo.

        \item Renda per capita por domicílio com TV a cabo (em USD).

        \item Custo médio mensal de manutenção de TV a cabo (em USD).

\end{itemize}
  
\end{itemize}

\end{frame}

% -----------------------------------------------------------------

\begin{frame}

\frametitle{Variáveis}

\begin{itemize}
  
  \itemsep 2ex
  
  \item Variáveis resposta:
  
    \begin{enumerate}
      \item Número de assinantes de TV a cabo (em milhares).
      \item Percentual de domicílios com TV a cabo.
    \end{enumerate}
    
  \item Variáveis explicativas:
  
    \begin{enumerate}
      \item Renda per capita por domicílio com TV a cabo (em USD).
      \item Custo médio mensal de manutenção de TV a cabo (em USD).
    \end{enumerate}
  
  \item No problema existem duas variáveis resposta, de diferentes naturezas: uma contagem e uma proporção.
  
  \item A estrutura do problema é ideal para utilizar um modelo multivariado.
  
\end{itemize}

\end{frame}

% -----------------------------------------------------------------

\begin{frame}

Para análise do problema, considere um modelo bivariado:

$$
g_r(\mu_r) = \beta_{r0} + \beta_{r1} renda + \beta_{r2} custo + \beta_{r3} renda:custo.
$$

Em que:

\begin{itemize}
    \itemsep 2ex

  \item O índice $r$ denota a variável resposta, r = 1,2.
    \begin{itemize}
      \item 1: Número de assinantes de TV a cabo (em milhares).
      \item 2: Percentual de domicílios com TV a cabo.
    \end{itemize}

  \item $\beta_{r0}$ denota o intercepto de cada resposta. 

  \item Temos três parâmetros de regressão para cada resposta:
  
  \begin{enumerate}
    
    \item $\beta_{r1}$ é o efeito de renda. 
    
    \item $\beta_{r2}$ é o efeito de custo.
    
    \item $\beta_{r3}$ representa o efeito da interação entre as variáveis renda e custo.
    
  \end{enumerate}

\item Considere ainda que:

  \begin{itemize}
    \item Cada resposta possui apenas um parâmetro de dispersão: $\tau_{r0}$.
    \item As unidades em estudo são independentes, logo $Z_0$ = \textbf{I}.
    \item Os parâmetros de potência foram fixados.
  \end{itemize}

\end{itemize}

\end{frame}

% -----------------------------------------------------------------

\begin{frame}

\frametitle{Pergunta}

Existe efeito de custo médio mensal de manutenção de TV a cabo (em USD) sobre o número de assinantes de TV a cabo (em milhares) \textbf{E} percentual de domicílios com TV a cabo?

$$H_0: \beta_{r2} = 0 \ vs \ H_1: \beta_{r2} \neq 0.$$ 

Ou, da mesma forma:

$$H_0: 
\begin{pmatrix}
\beta_{12} \\ 
\beta_{22}
\end{pmatrix} 
= 
\begin{pmatrix}
0 \\ 
0
\end{pmatrix}
\ vs \ 
H_1: 
\begin{pmatrix}
\beta_{12} \\ 
\beta_{22}
\end{pmatrix} 
\neq
\begin{pmatrix}
0 \\ 
0 
\end{pmatrix}.$$

\end{frame}

% -----------------------------------------------------------------

\begin{frame}

\frametitle{Efeito de custo com manutenção sobre número de assinantes\\ e percentual de domicílios com TV a cabo}

A hipótese pode ser reescrita na seguinte notação:

$$H_0: \boldsymbol{L}\boldsymbol{\theta_{\beta,\tau,p}} = \boldsymbol{c} \ vs \ H_1: \boldsymbol{L}\boldsymbol{\theta_{\beta,\tau,p}} \neq \boldsymbol{c}.$$ 

Em que:

\begin{itemize}
  
 \item $\boldsymbol{\theta_{\beta,\tau,p}^T}$ = $\begin{bmatrix} \beta_{10} \  \beta_{11} \ \beta_{12} \ \beta_{13} \ \beta_{20} \ \beta_{21} \ \beta_{22} \ \ \beta_{23} \ \tau_{10} \ \tau_{20} \end{bmatrix}$.


\item $\boldsymbol{L} = \begin{bmatrix} 0 & 0 & 1 & 0 & 0 & 0 & 0 & 0 & 0 & 0 \\
0 & 0 & 0 & 0 & 0 & 0 & 1 & 0 & 0 & 0 \end{bmatrix}$.
 
\item $\boldsymbol{c} = \begin{bmatrix} 0 \\ 0 \end{bmatrix}$, é o valor da hipótese nula. 

\end{itemize}

\end{frame}

% -----------------------------------------------------------------

\subsection{ANOVA \& MANOVA via teste Wald}

% ----------------------------------------------------------  -------

\begin{frame}[c, allowframebreaks]

\begin{center}

  {\normalsize \href{https://lineu96.github.io/st/}{ANOVA \& MANOVA via teste Wald}}
  
\end{center}

\end{frame}

% -----------------------------------------------------------------

\begin{frame}

\frametitle{ANOVA \& MANOVA via teste Wald}

\begin{itemize}
    \itemsep 2ex

  \item Com base na adaptação do teste Wald propostas, buscamos propor ANOVAs e MANOVAs via teste Wald.

  \item Propomos 3 tipos diferentes de análises de variância, nomeadas tipo I, II e III.

  \item Cada linha do quadro corresponde uma hipótese. Portanto, basta especificar uma matriz $\boldsymbol{L}$.

  \item Os procedimentos para análise de variância retornam um quadro para cada resposta.

  \item Os procedimentos para análise variância multivariadas retornam um único quadro. 
\end{itemize}

\end{frame}

% -----------------------------------------------------------------

\begin{frame}

\frametitle{ANOVA \& MANOVA tipo II}

\begin{itemize}
  \item São feitos testes comparando o modelo completo contra o modelo sem determinada variável. 
  
  \item Portanto, considerando o exemplo:

\begin{enumerate}
 \itemsep 2ex
  
  \item Testa se os parâmetros referentes a \textbf{renda} são iguais a 0. Ou seja, é avaliado o impacto da retirada de renda do modelo.
  
  \item Testa se os parâmetros referentes a \textbf{custo} são iguais a 0. Ou seja, é avaliado o impacto da retirada de custo do modelo.
  
  \item Testa se o efeito de \textbf{interação} é 0.

\end{enumerate}
\end{itemize}

\end{frame}

% -----------------------------------------------------------------

\subsection{Funções implementadas}

% -----------------------------------------------------------------

\begin{frame}[c, allowframebreaks]

\begin{center}

  {\normalsize \href{https://lineu96.github.io/st/}{Funções implementadas}}
  
\end{center}

\end{frame}

% -----------------------------------------------------------------

\begin{frame}
  \frametitle{Funções implementadas}

Baseando-nos nas funcionalidades do pacote \emph{car} \cite{car} e usando nossa adaptação do teste Wald implementamos uma série de funções:


\begin{table}[h]
\centering
\begin{tabular}{ll}
\hline
Função                   & Descrição \\ 
\hline

mc\_linear\_hypothesis() & Hipóteses lineares gerais especificadas pelo usuário \\

mc\_anova\_I()           & ANOVA  tipo I \\
mc\_anova\_II()          & ANOVA  tipo II \\
mc\_anova\_III()         & ANOVA  tipo III \\

mc\_manova\_I()          & MANOVA tipo I \\
mc\_manova\_II()         & MANOVA tipo II \\
mc\_manova\_III()        & MANOVA tipo III \\

mc\_anova\_disp()        & ANOVA  tipo III para dispersão \\
mc\_manova\_disp()       & MANOVA tipo III para dispersão \\

\hline
\end{tabular}
\caption{Funções implementadas}
\label{tab:funcoes}
\end{table}

\end{frame}

%-----------------------------------------------------------------------

\section{Próximas etapas}

% -----------------------------------------------------------------

\begin{frame}
  \frametitle{Próximas etapas}
 
\begin{itemize}
\itemsep 2ex

  \item Propor e implementar procedimentos para realização de testes de \textbf{comparações múltiplas}.
  
  \item Adequar os testes para que sejam válidos para diferentes \textbf{contrastes}. 
  
  \item Avaliar as propriedades e comportamento dos testes propostos com base em \textbf{estudos de simulação}.
  
  \item Motivar o potencial de aplicação das metodologias discutidas com base na \textbf{aplicação a conjuntos de dados reais}.
  
\end{itemize}

\end{frame}

% -----------------------------------------------------------------

\section{Considerações finais}

% -----------------------------------------------------------------

\begin{frame}
  \frametitle{Considerações finais}
  
\begin{itemize}
\itemsep 2ex
  \item O McGLM contorna importantes restrições encontradas nas classes clássicas de modelos.

    \item Nossa contribuição vai no sentido de fornecer ferramentas para uma melhor interpretação dos parâmetros estimados na classe.
    
    \item Nossa contribuição visa formas de responder sobre:
    
    \begin{enumerate}
      \item Importância das variáveis explicativas no problema.
      \item Impacto das medidas correlacionadas no conjunto de dados.
      \item Qual distribuição se adequa ao problema.
    \end{enumerate}

\end{itemize}

\end{frame}

%-----------------------------------------------------------------------

\begin{frame}[c, allowframebreaks]
  \frametitle{Referências bibliográficas}
  
  \begin{tiny}
    \bibliography{references}
  \end{tiny}
  
\end{frame}

\begin{frame}[c, allowframebreaks]

\begin{center}

  {\huge \href{https://lineu96.github.io/st/}{Obrigado!}}
  
  \vspace{0.5cm}
    
  {\normalsize \href{https://lineu96.github.io/st/}{Lineu Alberto Cavazani de Freitas}}
  
  {\normalsize \href{https://lineu96.github.io/st/}{lineuacf@gmail.com}}
  
  {\normalsize \href{https://lineu96.github.io/st/}{https://lineu96.github.io/st/}}
  
  {\normalsize \href{http://www.prppg.ufpr.br/ppginformatica/?lang=pb}{PPG Informática}}

\end{center}

\begin{center}
  \includegraphics[height=1.8cm]{img/capes_tp2.png}\hspace{2em}
  \includegraphics[height=1.8cm]{img/ufpr-transparent.png}\hspace{2em}
  \includegraphics[height=1.8cm]{img/dsbd-2x2-trans.png}
\end{center}

\end{frame}

%--------------------------------------------------

\end{document}
